\documentclass[a4paper,online]{adcv}
\usepackage[french]{babel}
\usepackage{booktabs}
\usepackage{graphicx}
\usepackage{multirow}
\usepackage[export]{adjustbox}
\usepackage{hyperref}

\hypersetup{
  colorlinks=true,
  linkcolor=E07A5F,
  filecolor=E07A5F,
  urlcolor=E07A5F,
  pdftitle={CV de Titouan Renard},
  pdfpagemode=FullScreen,
}

\title{CV de Titouan Renard}

\addbibresource{bib_titouan_renard.bib}

\adcvname{Titouan}{Renard}{}
\adcvtitle{Étudiant en économie politique}
\adcvaddress{}{}{Chemin de Mont-Riant 7, 2000 Neuchâtel}{Suisse}
\adcvemail{titouan.renard}{gmail}{com}
\adcvphone{(+41) 79 944 61 65}
\adcvwebsite{https://titouanrenard.eu/}{site web}
\adcvdate{Septembre 2026}

\begin{document}
Économiste de formation, issu d'un parcours d'ingénieur. Je m'intéresse principalement à l'économie internationale et à la politique commerciale, ainsi qu'à la vulgarisation et à la formation politique, que je pratique depuis plusieurs années dans le cadre associatif et médiatique. C'est ce goût pour la transmission qui me porte aujourd'hui vers des fonctions d'assistant parlementaire ou de chargé de communication.

\section{Formation}

\begin{adcvtabletwo}
  \adcvrowtwo{\textbf{Master en économie politique du capitalisme},
  Genève, Suisse}{2024-2026}
  \adcvrowmulti{Université de Genève (UNIGE)}
  \adcvrowtwo{\textbf{Master en économie},
  Barcelone, Espagne (semestre d'échange)}{2025}
  \adcvrowmulti{Universitat de Barcelona}
  \adcvrowtwo{\textbf{Master en robotique avec un mineur en science
  des données}, Lausanne, Suisse}{2023}
  \adcvrowmulti{EPFL, facultés STI et IC}
  \adcvrowtwo{\textbf{Bachelor en microtechnique}, Lausanne,
  Suisse}{2020}
  \adcvrowmulti{EPFL, faculté STI}
  \adcvrowtwo{\textbf{Maturité bilingue}, Neuchâtel, Suisse}{2016}
  \adcvrowmulti{Lycée bilingue français-anglais, Lycée Denis de
  Rougemont, option Physique et mathématiques appliquées}
\end{adcvtabletwo}



\section{Engagement politique, médiatique et associatif}

\begin{adcvtabletwo}
  \adcvrowtwo{\textbf{Politique} (Jeunes POP, mouvement de jeunesse du
  Parti Suisse du Travail)}{}
  \adcvrowtwo{\textit{Trésorier}}{2026}
  \adcvrowtwo{\textit{Coordinateur des sections francophones} (interrompue pendant mon échange en Espagne)}{2025}
  \adcvrowtwo{\textit{Responsable de la communication}}{2023-2024}
  \adcvrowtwo{\textbf{Médias}}{}
  \adcvrowtwo{\textit{Journaliste}, \href{https://www.youtube.com/@cestreelmedia/videos}{C'est Réel}, média en
  ligne}{2024-}
  \adcvrowmulti{Écriture, montage, prise de parole dans les vidéos du
    média et comptabilité.}
  \adcvrowtwo{\textit{Journaliste}, \href{https://voixpopulaire.ch/}{Voix Populaire}, journal de gauche radical Suisse}{2023-}
  \adcvrowmulti{Contributions occasionelles comme journaliste.}
  \adcvrowtwo{\textit{Journaliste}, \href{https://journal.unipoly.ch/}{Le Canard Huppé}, journal estudiantin Lausannois, papier et en ligne}{2020-2023}
  \adcvrowmulti{Écriture, travail d'édition, conception et gestion du site internet.}
  \adcvrowtwo{\textbf{Associatif}}{}
  \adcvrowtwo{\textit{Unipoly},
  association étudiante pour la durabilité à l'EPFL}{2019-2022}
  \adcvrowmulti{Membre du comité, puis président de l'association.}
  \adcvrowtwo{\textit{European Youth Parliament}, Suisse}{2015-2016}
  \adcvrowmulti{Expérience de prise de parole en public en anglais, première familiarisation avec les institutions de l'UE.}
\end{adcvtabletwo}

\section{Expérience}

\begin{adcvtabletwo}
  \adcvrowtwo{\textbf{Stages}}{}
  \adcvrowtwo{\textit{Stagiaire de recherche} au SYCAMORE (Systems
    Control and Multiagent Optimization Research) de l'EPFL, sous la
  supervision de Maryam Kamgarpour}{2023}
  \adcvrowmulti{Travail sur la complexité d'échantillonnage de
  l'apprentissage par renforcement inverse (travail publié).}
  \adcvrowtwo{\textit{Conception de projet de cours} au Laboratoire de
    neurosciences computationnelles, sous la supervision d'Alireza
  Modirshanechi}{2023}
  \adcvrowmulti{Conception du projet du cours « Artificial Neural
  Networks and Reinforcement Learning » à l'EPFL.}
  \adcvrowtwo{\textit{Stagiaire en apprentissage automatique} chez 3HLE
  Robotics and Automation}{2022-2023}
  \adcvrowmulti{Travail sur des modèles d'apprentissage profond pour
    l'\textit{estimation de pose 6D} à partir de données de capteurs
    RGB-D dans le contexte de la manipulation robotique industrielle,
  sous la supervision de Laurent Winkler}
  \adcvrowtwo{\textit{Stage d'été} au laboratoire de biorobotique de
  l'EPFL, sous la supervision d'Alessandro Crespi.}{2018}
  \adcvrowmulti{Implémentation d'un système simple de localisation par
  vision en C++ et OpenCV.}
  \adcvrowtwo{\textbf{Assistant étudiant à l'EPFL}}{}
  \adcvrowtwo{\textit{Supervision des séances d'exercices pour les de l'école étudiants}}{2018-2023}
  \adcvrowmulti{En physique, en informatique et en statistiques.}
  \adcvrowtwo{\textbf{Autre}}{}
  \adcvrowtwo{\textit{Webmaster} pour l'Institut de psychanalyse
  Charles Baudouin}{2025-2026}
  \adcvrowmulti{Développement, maintenance et mise à jour du site web
  de l'institut.}
\end{adcvtabletwo}


\section{Compétences}

\begin{adcvtabletwo}
  \adcvrowskip
  \adcvrowtwo{\textbf{Outils créatifs et numériques}, suite Adobe
  (Premiere Pro, Photoshop, After Effects), DaVinci Resolve, Figma,
  Canva ; création de sites web de zéro ou avec WordPress, community management et gestion des publications avec des outiles comme buffer.}
  \adcvrowskip
  \adcvrowtwo{\textbf{Compétences transversales}, prise de parole en
  public, écriture et gestion d'équipe acquises dans le cadre
  d'associations étudiantes et d'engagements politiques. Compétences d'ingénieur, mathématiques et programmation.}
  \adcvrowskip
  \adcvrowtwo{\textbf{Culture générale}, concernant l'économie internationale, la politique française, suisse et espagnole, et l'histoire des mouvements sociaux Européens.}{}
  \adcvrowskip
  \adcvrowtwo{\textbf{Langues}, français (langue maternelle), anglais
  (courant), allemand (B2), espagnol (B1)}{}
  \adcvrowskip
\end{adcvtabletwo}


\section{Projets}

\begin{adcvtabletwo}
  \adcvrowtwo{\textbf{Beyond Imbalances, A Competition Theoretic Account of the US-China Rivalry}, travail de master dans le cadre du Master in the Political Economy of Capitalism à l'université de Genève, sous la supervision du Prof. Cedric Durand et du Dr. Benjamin Burbaumer. Projet noté \textit{6/6}}{2026}
  \adcvrowmulti{Analyse hétérodoxe (Marxienne/Post-Keynesienne) de l'impact du second choc Chinois sur l'économie étatsunienne.}
  \adcvrowtwo{\textbf{Measuring Global Value Chains Integration Using Embodied Labour}, travail de séminaire pour
  le cours « Global Sustainability » à l'Universitat de Barcelona,
  sous la supervision de la Prof. Mònica Serrano. Projet noté
  \textit{10/10}}{2025}
  \adcvrowmulti{Analyse de la contribution du travail à la demande
  finale des différents secteurs et pays à l'aide de tableaux
  entrées-sorties interpays (ICIO). \href{https://titouanrenard.eu/documents/TRENARD2025_MeasuringGVCs.pdf}{Rapport du projet en ligne.}}
  \adcvrowtwo{\textbf{Swiss Financial Diplomacy in the Interwar Period, Preserving Bondholder’s Interests in the Context of the Renegotiation of the Ottoman Public Debt (1919-1928)}, projet de séminaire à l'UNIGE sous la
  supervision du Prof. Juan Flores Zendejas. Projet noté
  \textit{6/6}}{2025}
  \adcvrowmulti{Analyse de la collaboration entre les diplomates
  suisses et les porteurs d'obligations à partir de sources
  archivistiques diplomatiques. \href{https://titouanrenard.eu/documents/TRENARD2025_ottoman_financial_diplomacy.pdf}{Rapport du projet en ligne.}}
  \adcvrowtwo{\textbf{Provable convergence guarantees for constrained inverse reinforcement learning}, projet de
  master au SYCAMORE (Systems Control and Multiagent Optimization
  Research) de l'EPFL, sous la supervision de la Prof. Maryam
  Kamgarpour. Projet noté \textit{6/6}}{2023}
  \adcvrowmulti{Démonstration de garanties de convergence pour un
  algorithme d'apprentissage par renforcement inverse (IRL) sous
  contraintes ; le travail a conduit à une publication à l'IEEE CDC.}
\end{adcvtabletwo}

\section{Publications}

\begin{refsection}
  \nocite{morgan2025}
  \nocite{renardIRL}
  \nocite{desobeirPourLaTerre2021}
  \nocite{tumultePostCorona2020}
    \printbibliography[heading=none]
\end{refsection}

\end{document}
